\chapter{Những Quy Định Trên Trời}
\markboth{Những Quy Định Trên Trời}{Những Quy Định Trên Trời}

\section{Nội Quy Không Giải Thích}
\begin{truyen}{Nội Quy Không Giải Thích}{Rules Without Explanation}
\chuhoa{H}{ọc sinh:} ``Sao trường mình nhiều nội quy lạ vậy thầy? Cấm thì nhiều mà giải thích thì ít.''
\textit{(Student: ``Why does our school have so many strange rules? There are many prohibitions, but very few explanations.'')}

\teacherpair{Không phải nội quy nào cũng vô lý. Nhưng một quy định muốn được tôn trọng thì nên cho người ta thấy nó bảo vệ điều gì. Nếu chỉ quen ra lệnh mà không quen giải thích, nhà trường rất dễ dạy học sinh tuân thủ vì sợ chứ không phải vì hiểu.}{``Not every rule is unreasonable. But if a rule wants respect, it should show what it protects. If schools only issue commands and never explain, they teach students to obey out of fear rather than understanding.''}
\end{truyen}

\section{Cấm Cái Dễ Thấy, Bỏ Cái Cốt Lõi}
\begin{truyen}{Cấm Cái Dễ Thấy, Bỏ Cái Cốt Lõi}{Banning What Is Visible While Ignoring What Matters}
\chuhoa{H}{ọc sinh:} ``Có lúc em thấy người lớn soi tóc tai, giày dép kỹ hơn cả cách tụi em cư xử với nhau.''
\textit{(Student: ``Sometimes adults inspect hair and shoes more carefully than how we treat one another.'')}

\teacherpair{Đó là một phê bình đau nhưng đúng ở nhiều nơi. Cái gì dễ thấy thì dễ phạt. Còn những thứ khó hơn như thái độ, sự trung thực, cách nói năng thì cần công hơn mới uốn được. Khi hệ thống mê cái bề ngoài dễ quản, nó rất dễ bỏ quên phần nhân cách khó xây.}{``That criticism is painful but true in many places. What is visible is easy to punish. Harder things like attitude, honesty, and speech require more effort to shape. When a system becomes addicted to what is easy to manage, it forgets the character traits that are harder but more important to build.''}
\end{truyen}

\section{Nội Quy Tốt Là Nội Quy Nuôi Ý Thức}
\begin{truyen}{Nội Quy Tốt Là Nội Quy Nuôi Ý Thức}{Good Rules Build Awareness}
\chuhoa{H}{ọc sinh:} ``Vậy nội quy tốt là sao?''
\textit{(Student: ``Then what makes a good rule?'')}

\teacherpair{Là thứ vừa rõ ràng vừa có lý do. Nó giúp học sinh dần hiểu mình đang giữ cái gì, chứ không chỉ nơm nớp tránh bị bắt. Nội quy tốt không làm lớp học ngột ngạt hơn. Nó làm môi trường học bớt hỗn hơn.}{``A good rule is clear and has a reason. It helps students understand what they are preserving, not merely tremble to avoid punishment. Good rules do not make the classroom more suffocating. They make the learning environment less chaotic.''}
\end{truyen}

\begin{baihoc}
Quản học sinh bằng nội quy là cần.
\textit{(Rules are necessary in managing students.)}

Nhưng nội quy càng ít não, người tuân thủ càng ít tâm phục.
\textit{(But the less thought there is in a rule, the less sincere the compliance will be.)}
\end{baihoc}

\begin{ghinhoanh}
\textbf{Short English to remember:}

Rules need reasons.

Fear is weaker than understanding.
\end{ghinhoanh}

\begin{tuhocnhanh}
1. Em phản ứng với nội quy vì nó vô lý, hay vì em ghét bị giới hạn? \textit{(Do you react to rules because they are unreasonable, or because you hate limits?)}

2. Nếu được viết lại một nội quy, em sẽ sửa nó theo hướng nào? \textit{(If you could rewrite one rule, how would you improve it?)}
\end{tuhocnhanh}

\ngancach
